% Options for packages loaded elsewhere
\PassOptionsToPackage{unicode}{hyperref}
\PassOptionsToPackage{hyphens}{url}
\documentclass[
  11pt,
]{article}
\usepackage{xcolor}
\usepackage[left=3cm,right=3cm,top=2.5cm,bottom=2.5cm]{geometry}
\usepackage{amsmath,amssymb}
\setcounter{secnumdepth}{5}
\usepackage{iftex}
\ifPDFTeX
  \usepackage[T1]{fontenc}
  \usepackage[utf8]{inputenc}
  \usepackage{textcomp} % provide euro and other symbols
\else % if luatex or xetex
  \usepackage{unicode-math} % this also loads fontspec
  \defaultfontfeatures{Scale=MatchLowercase}
  \defaultfontfeatures[\rmfamily]{Ligatures=TeX,Scale=1}
\fi
\usepackage{lmodern}
\ifPDFTeX\else
  % xetex/luatex font selection
\fi
% Use upquote if available, for straight quotes in verbatim environments
\IfFileExists{upquote.sty}{\usepackage{upquote}}{}
\IfFileExists{microtype.sty}{% use microtype if available
  \usepackage[]{microtype}
  \UseMicrotypeSet[protrusion]{basicmath} % disable protrusion for tt fonts
}{}
\usepackage{setspace}
\makeatletter
\@ifundefined{KOMAClassName}{% if non-KOMA class
  \IfFileExists{parskip.sty}{%
    \usepackage{parskip}
  }{% else
    \setlength{\parindent}{0pt}
    \setlength{\parskip}{6pt plus 2pt minus 1pt}}
}{% if KOMA class
  \KOMAoptions{parskip=half}}
\makeatother
\usepackage{longtable,booktabs,array}
\usepackage{caption}
\captionsetup[table]{skip=6pt}
\newcounter{none} % for unnumbered tables
\usepackage{calc} % for calculating minipage widths
% Correct order of tables after \paragraph or \subparagraph
\usepackage{etoolbox}
\makeatletter
\patchcmd\longtable{\par}{\if@noskipsec\mbox{}\fi\par}{}{}
\makeatother
% Allow footnotes in longtable head/foot
\IfFileExists{footnotehyper.sty}{\usepackage{footnotehyper}}{\usepackage{footnote}}
\makesavenoteenv{longtable}
\usepackage{graphicx}
\makeatletter
\newsavebox\pandoc@box
\newcommand*\pandocbounded[1]{% scales image to fit in text height/width
  \sbox\pandoc@box{#1}%
  \Gscale@div\@tempa{\textheight}{\dimexpr\ht\pandoc@box+\dp\pandoc@box\relax}%
  \Gscale@div\@tempb{\linewidth}{\wd\pandoc@box}%
  \ifdim\@tempb\p@<\@tempa\p@\let\@tempa\@tempb\fi% select the smaller of both
  \ifdim\@tempa\p@<\p@\scalebox{\@tempa}{\usebox\pandoc@box}%
  \else\usebox{\pandoc@box}%
  \fi%
}
% Set default figure placement to htbp
\def\fps@figure{htbp}
\makeatother
% definitions for citeproc citations
\NewDocumentCommand\citeproctext{}{}
\NewDocumentCommand\citeproc{mm}{%
  \begingroup\def\citeproctext{#2}\cite{#1}\endgroup}
\makeatletter
 % allow citations to break across lines
 \let\@cite@ofmt\@firstofone
 % avoid brackets around text for \cite:
 \def\@biblabel#1{}
 \def\@cite#1#2{{#1\if@tempswa , #2\fi}}
\makeatother
\newlength{\cslhangindent}
\setlength{\cslhangindent}{1.5em}
\newlength{\csllabelwidth}
\setlength{\csllabelwidth}{3em}
\newenvironment{CSLReferences}[2] % #1 hanging-indent, #2 entry-spacing
 {\begin{list}{}{%
  \setlength{\itemindent}{0pt}
  \setlength{\leftmargin}{0pt}
  \setlength{\parsep}{0pt}
  % turn on hanging indent if param 1 is 1
  \ifodd #1
   \setlength{\leftmargin}{\cslhangindent}
   \setlength{\itemindent}{-1\cslhangindent}
  \fi
  % set entry spacing
  \setlength{\itemsep}{#2\baselineskip}}}
 {\end{list}}
\usepackage{calc}
\newcommand{\CSLBlock}[1]{\hfill\break\parbox[t]{\linewidth}{\strut\ignorespaces#1\strut}}
\newcommand{\CSLLeftMargin}[1]{\parbox[t]{\csllabelwidth}{\strut#1\strut}}
\newcommand{\CSLRightInline}[1]{\parbox[t]{\linewidth - \csllabelwidth}{\strut#1\strut}}
\newcommand{\CSLIndent}[1]{\hspace{\cslhangindent}#1}
\setlength{\emergencystretch}{3em} % prevent overfull lines
\providecommand{\tightlist}{%
  \setlength{\itemsep}{0pt}\setlength{\parskip}{0pt}}
% Statistics in Medicine submission preamble.
% Shared across analysis/report/<paper>/report.Rmd files via
%   includes:
%     in_header: ../sim-preamble.tex
% in the YAML output block.
%
% The page footer (source path, render time, version) is no
% longer set here. It is supplied at render time by the vendored
% tools/stamp-render.R, which injects tools/stamp.tex.

% Continuous line numbering for editorial review. The 'mathlines'
% option extends numbering through display-math environments
% (equation, align, ...), which SIM editorial workflows expect.
\usepackage[mathlines]{lineno}
\linenumbers
\usepackage{amsmath}

% Spacing: linestretch is set in YAML (1.5 line spacing).
\usepackage{setspace}
%% tools/stamp.tex
%% Provenance footer for PDFs rendered from this project.
%% Vendored by zzcollab; this file is not project-specific. The
%% footer prints the source document and the time of rendering.
%% The git version is defined here too, and so is available to a
%% project that wants it on the page, but it is left off the
%% default footer: it already names the staged copy in share/ and
%% appears in its manifest row. All values are supplied at render
%% time by tools/stamp-render.R, which writes a small generated
%% file that redefines the macros below. The \providecommand
%% defaults keep a bare LaTeX compile working when the document is
%% built without the wrapper.
\usepackage{fancyhdr}
\usepackage{xcolor}
\providecommand{\stampsource}{\detokenize{(source not recorded)}}
\providecommand{\stamptime}{(time not recorded)}
\providecommand{\stampversion}{\detokenize{(version not recorded)}}
\setlength{\footskip}{40pt}
\newcommand{\stampfootcontent}{%
  \footnotesize\color{gray}%
  \parbox[b]{\textwidth}{%
    \stamptime\hfill\thepage\\[2pt]%
    \ttfamily\raggedright\stampsource}}
\pagestyle{fancy}
\fancyhf{}
\fancyfoot[C]{\stampfootcontent}
\renewcommand{\headrulewidth}{0pt}
\renewcommand{\footrulewidth}{0.4pt}
%% The title page uses the `plain` page style; redefine it so the
%% stamp appears there too.
\fancypagestyle{plain}{%
  \fancyhf{}%
  \fancyfoot[C]{\stampfootcontent}%
  \renewcommand{\headrulewidth}{0pt}%
  \renewcommand{\footrulewidth}{0.4pt}}
\renewcommand{\stampsource}{\textasciitilde{}/\allowbreak{}prj/\allowbreak{}res/\allowbreak{}36-pmsimstats-ng/\allowbreak{}pmsimstats-ng/\allowbreak{}analysis/\allowbreak{}report/\allowbreak{}13-two-component-dgp/\allowbreak{}report.Rmd}
\renewcommand{\stamptime}{2026-09-08 19:00 PDT}
\renewcommand{\stampversion}{\detokenize{955c604-wip}}
\usepackage{amsmath}
\usepackage{amssymb}
\usepackage{booktabs}
\usepackage{array}
\usepackage{longtable}
\usepackage{float}
\usepackage[mathlines]{lineno}
\linenumbers
\usepackage{bookmark}
\IfFileExists{xurl.sty}{\usepackage{xurl}}{} % add URL line breaks if available
\urlstyle{same}
% fallback for those not using the hyperref driver hyperxmp:
\makeatletter
\@ifundefined{xmpquote}{\newcommand{\xmpquote}[1]{#1}}{}
\makeatother
\hypersetup{
  pdftitle={Data generation machinery for N-of-1 clinical trial simulations under a two-component response decomposition},
  pdfauthor={pmsimstats team},
  hidelinks,
  pdfcreator={LaTeX via pandoc}}

\title{Data generation machinery for N-of-1 clinical trial simulations under a two-component response decomposition}
\author{pmsimstats team}
\date{September 08, 2026}

\begin{document}
\maketitle

{
\setcounter{tocdepth}{2}
\tableofcontents
}
\setstretch{1.1}
\section{Abstract}

\textbf{Background.} Simulation studies of aggregated N-of-1 trials
must choose where in the joint distribution a biomarker-treatment
interaction lives. A companion paper compares two such choices, mean
moderation and covariance moderation, under a three-component
decomposition of treatment response into pharmacological (BR),
placebo-belief (PB), and time-variant natural-history (TV)
components. The TV component carries no biomarker coupling in either
architecture, so its role in that comparison is not obvious.

\textbf{Methods.} We repeat the comparison under a reduced
data-generating process carrying only BR and PB. Dropping TV reduces
the correlation matrix from \((2 + 3n)\) to \((2 + 2n)\), from \(26 \times
26\) to \(18 \times 18\) at the eight occasions used throughout this
program. We cross two architectures against three trial designs, two
sample sizes, three biomarker-moderation levels, and three carryover
half-lives, holding the analysis model fixed. Appendix C derives the
positive-definiteness constraint on the moderation parameter in
closed form and characterizes what the reduction does and does not
relax.

\textbf{Results.} \emph{[PLACEHOLDER: pending the production run.
To report: the carryover-response profile of each architecture under
the reduced DGP; Type I error at the null; the comparison of absolute
power against the three-component setting.]}

\textbf{Conclusions.} \emph{[PLACEHOLDER: pending results. Expected
to state the conditions under which the reduction is appropriate for
methodological work and the conditions under which the full
decomposition remains necessary.]}

\section{1. Introduction}

An aggregated N-of-1 trial pools repeated within-participant on-drug
and off-drug contrasts to estimate a population-level
biomarker-treatment interaction\textsuperscript{\citeproc{ref-hendrickson2020}{1}}. Simulation
studies of such trials must specify a data-generating process, and
that specification involves two choices usually made implicitly. The
first is where the interaction lives in the joint distribution of
biomarker and response. The second is how many distinct latent
components the response trajectory is decomposed into.

The companion paper 01 takes up the first choice. It compares
\emph{mean moderation}, in which the biomarker shifts the mean of the
pharmacological response, against \emph{covariance moderation}, in
which the biomarker's correlation with that response is
treatment-state dependent. It shows the two are not interchangeable:
they respond differently to carryover, and the difference is large
enough to change a design recommendation.

The present paper takes up the second choice. The program's standard
DGP decomposes response into three components: BR, the
pharmacological action of the compound; PB, the improvement driven by
the participant's belief that treatment will work; and TV, drift in
the underlying condition independent of treatment. That decomposition
is developed and defended in paper 06, whose subject is the
inferential cost of failing to decompose.

There is a specific reason to ask whether all three are needed for
\emph{methodological} work. The biomarker couples to BR in both
architectures. It couples to PB only through an optional
contamination parameter that is zero by default. It does not couple
to TV at all: the package provides no parameter for a biomarker-TV
correlation, so that entry of the correlation matrix is identically
zero in every simulation this program has run. Paper 06's
omitted-variable identity then implies that TV cannot displace the
biomarker-treatment slope, and paper 06's Study A confirms this
empirically, finding no trend in bias across its TV magnitude axis.

If TV cannot bias the estimand, what does it do? It contributes
variance. It also contributes \(n\) rows and columns to the correlation
matrix, two cross-component blocks, and a set of
positive-definiteness constraints. Whether that cost is repaid in a
study whose object is the comparison between architectures is the
question this paper answers, by repeating paper 01's comparison
without it.

We hold everything else fixed. The trial designs, the analysis model,
the parameter calibration, the biomarker-moderation levels, and the
carryover half-lives are those of paper 01. Section 2 states the
model and the two architectures under the reduced decomposition,
Section 3 gives the results, and Section 4 takes up what follows for
the design of future simulation work in this program. Three
appendices carry the structural material: Appendix A on why the two
architectures are treated as a pair, Appendix B on the correlation
matrix the reduced DGP assembles, and Appendix C on the
positive-definiteness constraint the reduction was expected to
relax.

\section{2. Methods}

\subsection{2.1 Notation and analysis model}

Let \(i = 1, \dots, N\) index participants and \(p = 1, \dots, n\)
measurement occasions at cumulative weeks \(w_1 < \cdots < w_n\). Write
\(B_i\) for the participant's baseline biomarker, \(\mathrm{BL}_i\) for
baseline symptom severity, and \(S_{x,ip}\) for the observed symptom
score. Under the two-component decomposition,

\[
S_{x,ip} = \mathrm{BL}_i - \big( BR_{ip} + PB_{ip} \big),
\]

so that a beneficial response lowers the score. The three-component
model of paper 06 adds \(TV_{ip}\) inside the parentheses; setting it
to zero recovers the present model exactly.

Both components are drawn jointly across occasions from a
multivariate normal whose means follow modified Gompertz
trajectories,

\[
\mathrm{mg}(t; a, d, r) = a \exp\!\big(-d \exp(-r t)\big),
\]

evaluated at a component-specific time argument: cumulative time on
drug \(t_{od}\) for BR, and accumulated belief exposure \(t_{pb}\) scaled
by the design's expectancy weight \(\eta\) for PB. The distinct time
arguments are what make the components separable from design
contrasts.

The analysis model is held fixed across all conditions,

\[
S_{x,ip} = \beta_0 + \beta_{bm} B_i + \beta_t t_p
  + \beta_D D_{bc,ip} + \beta_{bm:D} B_i D_{bc,ip}
  + u_i + \varepsilon_{ip},
\]

fitted with \texttt{nlme::lme}, a participant random intercept \(u_i
\sim N(0, \sigma_u^2)\), and continuous-time AR(1) residual
correlation \texttt{corCAR1(form = \textasciitilde\ t | ptID)}. Here
\(D_{bc,ip}\) is the continuous exposure indicator, equal to one on
drug and decaying as \(\exp(-\lambda t_{sd,ip})\) off drug with
\(\lambda = \ln 2 / t_{1/2}\). The estimand is \(\beta_{bm:D}\).
Appendix A.4 states why the analysis model is held fixed.

\subsection{2.2 Two data-generating architectures}

\subsubsection{2.2.1 Mean moderation (Architecture A)}

Under mean moderation the biomarker shifts the conditional mean of
the pharmacological component at on-drug occasions. After the
multivariate normal draw, BR is displaced by

\[
BR_{ip} \leftarrow BR_{ip} + c_{bm}\, z_i\, \sigma_{BR}
  \quad \text{for occasion } p \text{ on drug},
\]

where \(z_i = (B_i - \mu_{bm})/\sigma_{bm}\) is the standardized
biomarker. The covariance structure is untouched: the biomarker is
independent of every component draw, and the interaction is carried
entirely by the first moment. In the notation of Appendix B, mean
moderation sets the coupling vector \(b \equiv 0\).

\subsubsection{2.2.2 Covariance moderation (Architecture B)}

Under covariance moderation the means are untouched and the
interaction is carried by a treatment-state-dependent correlation,

\[
\mathrm{Cor}(B_i, BR_{ip}) = c_{bm} \cdot \phi(t_{sd,ip}), \qquad
\phi(t_{sd}) = \begin{cases}
1, & \text{on drug} \\
\exp(-\lambda t_{sd}), & \text{off drug.}
\end{cases}
\]

The biomarker is more strongly correlated with the pharmacological
response while the drug is present, and that association decays after
discontinuation at the pharmacokinetic rate.

\subsubsection{2.2.3 What the reduced decomposition changes}

Neither architecture references TV. Removing it therefore leaves both
definitions above unchanged and alters the simulation in three
mechanical ways.

First, the correlation matrix shrinks, from \(2 + 3n\) to \(2 + 2n\).
Appendix B gives the resulting block partition.

Second, the outcome loses a variance term. At the reference
calibration the three components carry standard deviations of 10
(TV), 10 (PB), and 8 (BR), so TV accounts for approximately 38\% of
the component variance. Removing it reduces the residual variation
against which the interaction is tested, and absolute power rises
accordingly. This is an expected arithmetic consequence rather than a
finding, and it is why absolute power in this paper is not comparable
to paper 01's.

Third, the cross-component structure simplifies. The three-component
matrix carries three cross-component blocks; the two-component matrix
carries one. Appendix B.6 sets out exactly what is removed, and
Appendix C shows that what is removed is not what constrains the
moderation parameter.

\subsection{2.3 The carryover recursion}

Both architectures apply carryover to the BR mean trajectory: at an
off-drug occasion, a decayed fraction of the pre-discontinuation
response persists. The implementation used here differs from the one
behind paper 01's published results, and the difference is worth
stating.

The earlier implementation computed

\[
\mu_{BR,p} \leftarrow \mu_{BR,p} + \mu_{BR,p-1}\,
  (1/2)^{t_{sd,p}/t_{1/2}},
\]

in which \(\mu_{BR,p-1}\) is the \emph{already adjusted} value from the
preceding occasion while \(t_{sd,p}\) is \emph{cumulative} time since
discontinuation. Elapsed time is therefore counted twice, and every
off-drug occasion after the first in an uninterrupted run is
over-decayed. At a one-week half-life the second off-drug occasion
receives half the residual effect it should.

The present implementation anchors the decay to the mean at
discontinuation,

\[
\mu_{BR,p} \leftarrow \mu_{BR,p} + \mu_{BR,p^{\ast}}\,
  (1/2)^{t_{sd,p}/t_{1/2}},
\]

where \(p^{\ast}\) indexes the most recent on-drug occasion, so the
cumulative decay factor is applied once. We verified at three
reference cells (Hybrid, OL+BDC, and CO at \(N = 70\), \(c_{bm} = 0.45\),
\(t_{1/2} = 1.0\), 400 replicates under a common seed) that the
correction moves power by at most 0.005 against a Monte Carlo
standard error of about 0.025. It is adopted because it is correct,
not because it changes results.

\subsection{2.4 Simulation design}

The grid crosses architecture, trial design, sample size,
biomarker-moderation strength, and carryover half-life.

{\def\LTcaptype{none} % do not increment counter
\begin{longtable}[]{@{}ll@{}}
\toprule\noalign{}
Factor & Levels \\
\midrule\noalign{}
\endhead
\bottomrule\noalign{}
\endlastfoot
Architecture & mean moderation (A), covariance moderation (B) \\
Trial design & CO, Hybrid, OL+BDC \\
Sample size \(N\) & 35, 70 \\
Biomarker moderation \(c_{bm}\) & 0, 0.30, 0.45 \\
Carryover half-life \(t_{1/2}\) & 0, 0.5, 1.0 weeks \\
Components & BR + PB \\
Replicates & 500 \\
\end{longtable}
}

This gives \(2 \times 3 \times 2 \times 3 \times 3 = 108\) cells. The
\(c_{bm} = 0\) rows are the Type I error check against nominal 5\%; the
\(t_{1/2} = 0\) rows give a no-carryover reference against which the
cost of nonzero carryover is read off directly.

Throughout, \(N\) is the \textbf{total} number of participants,
allocated as evenly as possible across the design's randomization
paths, so a design with \(P\) paths assigns about \(N/P\) to each.

Absolute power is not comparable between architectures at the same
nominal \(c_{bm}\), since the parameter is not calibrated to equal
power across the two channels. The interpretable comparison is the
carryover profile within each architecture, and the qualitative shape
of that profile between them.

\section{3. Results}

\emph{[PLACEHOLDER pending the production run.]}

Intended structure, mirroring paper 01 Section 3:

\begin{itemize}
\item \textbf{3.1 Power under carryover at $N = 70$.} Table of power
  by architecture, design, and $t_{1/2}$ at $c_{bm} = 0.45$, with
  MCSEs and non-convergence rates. The question is whether the
  covariance arm's decline with $t_{1/2}$ and the mean arm's
  flatness, reported by paper 01 under three components, survive the
  reduction.
\item \textbf{3.2 Mechanism.} Whether the divergence is attributable
  to the two channels through which carryover attacks covariance
  moderation (mean blurring and correlation erosion) against the one
  channel available under mean moderation.
\item \textbf{3.3 Type I error and estimator bias.} Rejection rate at
  $c_{bm} = 0$ against nominal 0.05, across designs, sample sizes,
  and half-lives.
\item \textbf{3.4 Comparison with the three-component setting.}
  Absolute power against paper 01's published figures, with the
  variance argument of Section 2.2.3 as the expected explanation.
  Note that this comparison is qualitative: paper 01's figures were
  produced under the earlier carryover recursion and a partly
  different grid, so a precise decomposition of the difference into
  its variance and recursion parts is not available without a matched
  run.
\end{itemize}

\section{4. Discussion}

\emph{[PLACEHOLDER pending results.]}

Intended structure:

\begin{itemize}
\item \textbf{4.1 What the TV component is for.} The narrow claim
  the results can support (TV is inert for this estimand under the
  DGP as implemented) against the broad claim they cannot (that TV
  is dispensable). Three surviving roles: attribution estimands,
  biomarkers correlated with natural-history drift, and absolute
  power figures calibrated against a clinical reference.
\item \textbf{4.2 When the reduction is appropriate.} Proposed
  division: reduced DGP for comparative methodological work, full
  DGP whenever a number is read as an estimate of what a real trial
  would achieve or the decomposition is itself the object of study.
\item \textbf{4.3 A gap this paper exposes.} Paper 06 states its
  identity over both non-pharmacological components and recommends
  decomposition when a biomarker may correlate with PB \emph{or TV},
  but the package implements \texttt{c.bm.pb} and no TV counterpart.
  The TV half of that claim has never been exercised. Our own
  conclusion is correspondingly conditional: TV is inert under the
  DGP as implemented, not intrinsically.
\item \textbf{4.4 Limitations.} Single parameter calibration; removal
  rather than flattening (Section 2.2.3 and Appendix B.6); no matched
  three-component run.
\end{itemize}

\section{5. Conclusions}

\emph{[PLACEHOLDER pending results.]}

\section{Conflict of interest}

The authors declare no conflicts of interest.

\section{Funding}

This work received no specific funding.

\section{Reproducibility}

The simulation driver is
\texttt{analysis/scripts/two-component-dgp/01-run-grid.R} and its
output is versioned at
\texttt{analysis/data/13-two-component-grid.rds}. The driver sets a
master seed and derives a per-cell seed from it. Every number quoted
in the results is read from that file at render time rather than
typed in.

The two-component decomposition is supported by the
\texttt{components} argument added to \texttt{generateData()} and
\texttt{buildSigma()} in \texttt{R/generateData.R}. It defaults to
\texttt{c("tv", "pb", "br")}, so existing calls are unaffected;
passing \texttt{c("pb", "br")} produces the reduced matrix used here.
The argument requires that \texttt{"br"} be retained, since BR
carries the interaction channel in both architectures, and
canonicalizes the ordering so that label and standard-deviation
sequence is stable regardless of how the caller writes it. The
package test suite passes unchanged under the default.

The corrected carryover recursion of Section 2.3 is applied in
\texttt{R/generateData.R} and
\texttt{implementations/tidyverse/R/functions.R}. Other
implementation collections retain the earlier form; see
\texttt{02-deviations.md} in the driver directory.

\bigskip\noindent

\rule{\textwidth}{0.4pt}\bigskip

\section{Appendix A. Why mean and covariance moderation}

\subsection{A.1 The two channels exhaust the Gaussian case}

The simulations draw participant trajectories from a multivariate
normal distribution, and a multivariate normal is completely
determined by its mean vector \(\mu\) and its covariance matrix
\(\Sigma\). Any biomarker-treatment interaction written into such a
data-generating process must therefore enter through \(\mu\), through
\(\Sigma\), or through both. Mean moderation is the first route,
covariance moderation the second, and the combined architecture
developed in a companion paper\textsuperscript{\citeproc{ref-pmsimstats-paper11}{2}} is the third.
Within the Gaussian family there is no fourth place to put the
interaction. This is why the two architectures are treated as a pair
rather than as two entries in an open-ended list: they are not two
options among many, but the two components of a decomposition.

Reducing the response decomposition from three components to two does
not disturb this argument. The exhaustiveness claim concerns where an
interaction can be written into a Gaussian joint distribution, not
how many latent components that distribution is built from. Dropping
TV removes rows from \(\Sigma\); it does not create a third place for
the interaction to live, nor remove either of the two that exist.

The exhaustiveness claim is specific to the Gaussian family and
should not be read more broadly. A data-generating process built on a
finite mixture, or on explicit latent responder classes, can carry an
interaction in features that neither \(\mu\) nor \(\Sigma\) captures for
a single component, and Senn notes finite mixtures as an alternative
variance-component model\textsuperscript{\citeproc{ref-senn2019samplesize}{3}}. Covariance moderation
is best understood as the second-moment reduced form of such a
mixture, which is why the latent-class treatment is developed
separately\textsuperscript{\citeproc{ref-pmsimstats-paper03}{4}}.

\subsection{A.2 What each channel asserts}

The two channels are not merely different parameterizations of one
mechanism; each commits to a distinct claim about what the biomarker
is.

Mean moderation asserts that the biomarker fixes a participant's drug
effect. Two participants sharing a biomarker value share an expected
treatment effect exactly, and any difference between their realized
responses is measurement error. This is the appropriate commitment
when the biomarker measures pathway engagement directly, or stands in
a causal path between treatment and response.

Covariance moderation asserts that the biomarker predicts response
only probabilistically. Two participants sharing a biomarker value
may still respond differently, because the biomarker indexes an
unobserved responder status imperfectly rather than determining
response. This is the appropriate commitment when the biomarker is a
statistical proxy, which is the situation the prazosin and PTSD
application was designed around.

The distinction is substantive rather than technical. It determines
what a fitted interaction coefficient means, and how much power
carryover costs.

\subsection{A.3 Standing in the literature}

Mean moderation is the prevailing convention wherever
biomarker-treatment interactions are simulated. Haller and Ulm\textsuperscript{\citeproc{ref-haller2018}{5}} place the interaction in a regression mean, as do the
oncology design literatures on enrichment and predictive-marker
trials\textsuperscript{\citeproc{ref-simon2010}{6},\citeproc{ref-freidlin2014}{7}}. Grenet and colleagues\textsuperscript{\citeproc{ref-grenet2020}{8}} use the same form for a crossover. That literature is,
however, overwhelmingly parallel-group.

The N-of-1 methodological literature does not adopt mean moderation,
because it does not model an observed moderator at all. Its standard
formulation carries a random participant-by-treatment interaction, a
variance component describing unexplained heterogeneity with no
covariate attached\textsuperscript{\citeproc{ref-senn2019samplesize}{3},\citeproc{ref-senn2016mastering}{9},\citeproc{ref-wang2019}{10}}. Biomarker-validation simulation for N-of-1 designs is
therefore close to unoccupied territory rather than a settled field
with a convention to follow, and Hendrickson et al.\textsuperscript{\citeproc{ref-hendrickson2020}{1}} is its principal entry. The two-channel
decomposition of A.1 is adopted precisely because no established
N-of-1 convention exists to inherit.

\subsection{A.4 Holding the analysis model fixed}

Both architectures are analyzed with the identical model of Section
2.1. This is deliberate and worth stating explicitly, because a
data-generating process and an analysis model are separate objects,
and a comparison that varies both attributes differences to neither.
The continuous exposure variable \texttt{Dbc} loses contrast variance
as carryover grows, and a binary on-drug indicator does not. A
comparison in which one architecture carried a continuous regressor
and the other a binary one would confound that effect with the
architecture contrast it was meant to isolate. Holding the analysis
model fixed means every difference reported in Section 3 is
attributable to the data-generating process alone.

The same argument governs the reduction studied here. The analysis
model does not reference the component decomposition at all: it sees
only \(S_x\), and \(S_x\) is a sum. Removing TV therefore changes what
the analysis is fitted to without changing what it fits, which is the
property that makes the comparison in Section 3 interpretable.

\bigskip\noindent

\rule{\textwidth}{0.4pt}\bigskip

\section{Appendix B. Correlation structure under the two-component
decomposition}

This appendix sets out the correlation matrix the reduced
data-generating process assembles, in the notation of paper 01's
Appendix B so that the two can be read side by side. Symbols follow
the variable names used in the code: \(\mathrm{BM}\) is the biomarker,
written \(B_i\) in Section 2.1, and \(BR_p\), \(PB_p\) the response
components at occasion \(p\).

Mean moderation appears below only through the statement that its
coupling vector is zero. It writes no biomarker-response correlation
into the covariance at all; its interaction is carried in the mean
instead, as Section 2.2.1 sets out.

\subsection{B.1 Variable ordering and block partition}

Each participant's trajectory is drawn from a single multivariate
normal over an ordered variable list. With \(n\) measurement occasions
at cumulative weeks \(w_1 < \cdots < w_n\), the two-component vector is

\[
X = \bigl(\,\mathrm{BM},\; \mathrm{BL},\;
PB_1,\dots,PB_n,\;
BR_1,\dots,BR_n \,\bigr)^{\!\top},
\qquad \dim X = 2 + 2n ,
\]

against \(2 + 3n\) under the three-component decomposition. The
ordering is component-major: all \(n\) occasions of one component are
contiguous, rather than all components at one occasion. Writing \(A_c\)
for the \(n \times n\) within-component block of component \(c \in
\{pb, br\}\) and \(C_{pb,br}\) for the single cross-component block, the
correlation matrix partitions as

\[
R \;=\;
\begin{pmatrix}
1 & 0 & 0^{\top}   & b^{\top} \\[2pt]
0 & 1 & 0^{\top}   & 0^{\top} \\[2pt]
0 & 0 & A_{pb}     & C_{pb,br} \\[2pt]
b & 0 & C_{br,pb}  & A_{br}
\end{pmatrix}.
\]

Three features carry over unchanged from the three-component case.
The baseline term \(\mathrm{BL}\) is uncorrelated with every other
variable, so its row and column are those of the identity. The
biomarker \(\mathrm{BM}\) couples only to the response component \(BR\),
through the vector \(b \in \mathbb{R}^{n}\); its correlation with \(PB\)
is zero throughout this paper. And \(b\) is the only place in \(R\) where
the architectures differ.

The third feature deserves emphasis because it is what licenses the
reduction. In the three-component matrix the biomarker's row is
\((1, 0, 0^{\top}, 0^{\top}, b^{\top})\): zero against both \(TV\) and
\(PB\). Deleting the \(TV\) rows and columns therefore deletes only
zeros from that row. The coupling vector is untouched, entry for
entry.

\subsection{B.2 Within-component blocks $A_c$}

Let \(\rho_c\) denote the autocorrelation parameter of component \(c\)
(\texttt{c.pb}, \texttt{c.br} in the code). The correlation decays
with the elapsed time between occasions, using the cumulative week
gap rather than the index difference:

\[
[A_c]_{ij} = \rho_c^{\,|w_i - w_j|}.
\]

For \(n = 4\), writing \(d_{ij} = |w_i - w_j|\),

\[
A_c =
\begin{pmatrix}
1 & \rho_c^{d_{12}} & \rho_c^{d_{13}} & \rho_c^{d_{14}}\\
\rho_c^{d_{12}} & 1 & \rho_c^{d_{23}} & \rho_c^{d_{24}}\\
\rho_c^{d_{13}} & \rho_c^{d_{23}} & 1 & \rho_c^{d_{34}}\\
\rho_c^{d_{14}} & \rho_c^{d_{24}} & \rho_c^{d_{34}} & 1
\end{pmatrix}.
\]

Because the exponent uses calendar time, unequally spaced designs
produce unequal decay across consecutive occasions. In the Hybrid
design, whose occasions fall at weeks \(4, 8, 9, 10, 11, 12, 16, 20\),
consecutive gaps of one week during the discontinuation phase yield
much higher adjacent correlations than the four-week gaps at either
end.

We note that the code uses the cumulative week gap, not the occasion
index difference, and that this choice is consequential: Appendix C.4
reports that the index-based form would cost roughly a fifth of the
feasible range of \(c_{bm}\).

\subsection{B.3 The cross-component block $C_{pb,br}$}

Let \(c_1\) denote the same-occasion cross-component correlation
(\texttt{c.cf1t}) and \(c_{\times}\) the different-occasion
cross-component correlation (\texttt{c.cfct}). The same-occasion
entry is a constant and the different-occasion entries inherit the
AR(1) decay:

\[
[C_{pb,br}]_{ij} =
\begin{cases}
c_1, & i = j\\
c_{\times}\,\rho^{\,|w_i - w_j|}, & i \neq j .
\end{cases}
\]

Under the three-component decomposition there are three such blocks,
\(C_{tv,pb}\), \(C_{tv,br}\) and \(C_{pb,br}\), all of this form. The
reduction retains the last and deletes the first two.

\subsection{B.4 Biomarker coupling vector $b$}

The vector \(b\) carries the biomarker-treatment interaction under
covariance moderation. Its entry is gated on the drug indicator and
decays with time since discontinuation \(t_{sd,p}\) at rate \(\lambda =
\ln 2 / t_{1/2}\):

\[
b_p =
\begin{cases}
c_{bm}, & \text{occasion } p \text{ on drug}\\
c_{bm}\,e^{-\lambda t_{sd,p}}, & \text{occasion } p \text{ off drug} .
\end{cases}
\]

Under mean moderation \(b \equiv 0\) and the interaction is carried by
the displacement of \(\mu_{BR}\) described in Section 2.2.1. The two
architectures are therefore distinguished, in the language of this
appendix, by exactly one object: whether \(b\) is populated or zero.

\subsection{B.5 A worked example}

To make the blocks concrete, take four occasions drawn from the
Hybrid design at cumulative weeks \(w = (10, 11, 12, 16)\), with the
first and fourth on drug and the second and third off it, so that
\(t_{sd} = (0, 1, 2, 0)\). Set \(\rho_{br} = 0.6\), \(\rho_{pb} = 0.4\),
\(c_1 = 0.3\), \(c_{\times} = 0.1\), \(c_{bm} = 0.45\), and a carryover
half-life of one week, giving \(\lambda = \ln 2 \approx 0.693\). The
matrix of week gaps is

\[
\bigl(|w_i - w_j|\bigr) =
\begin{pmatrix}
0 & 1 & 2 & 6\\
1 & 0 & 1 & 5\\
2 & 1 & 0 & 4\\
6 & 5 & 4 & 0
\end{pmatrix}.
\]

The unequal spacing matters: occasions one to three are a week apart,
while the fourth sits four weeks after the third.

\subsubsection{B.5.1 Within-component block $A_{br}$}

With \(\rho_{br} = 0.6\),

\[
A_{br} =
\begin{pmatrix}
1 & 0.6000 & 0.3600 & 0.0467\\
0.6000 & 1 & 0.6000 & 0.0778\\
0.3600 & 0.6000 & 1 & 0.1296\\
0.0467 & 0.0778 & 0.1296 & 1
\end{pmatrix}.
\]

The fourth occasion is nearly uncorrelated with the first three, a
direct consequence of the four-week gap: \(0.6^{4} = 0.1296\) against
\(0.6^{1} = 0.6\) for an adjacent pair.

\subsubsection{B.5.2 Cross-component block $C_{pb,br}$}

With \(c_1 = 0.3\), \(c_{\times} = 0.1\) and the decay taken at
\(\rho = 0.6\),

\[
C_{pb,br} =
\begin{pmatrix}
0.3000 & 0.0600 & 0.0360 & 0.0047\\
0.0600 & 0.3000 & 0.0600 & 0.0078\\
0.0360 & 0.0600 & 0.3000 & 0.0130\\
0.0047 & 0.0078 & 0.0130 & 0.3000
\end{pmatrix}.
\]

The block is symmetric here because the two components share a decay
rate in this example; in general it need not be.

\subsubsection{B.5.3 Biomarker vector $b$}

Occasions one and four are on drug, occasions two and three off it
with \(t_{sd} = 1\) and \(2\). With \(c_{bm} = 0.45\) and \(\lambda = \ln 2\),

\[
b = \bigl(0.4500,\; 0.2250,\; 0.1125,\; 0.4500\bigr)^{\!\top}.
\]

The off-drug entries halve with each additional week since
discontinuation, and the fourth occasion returns to full strength
because the participant is back on drug. It is this alternation that
carries the interaction, and its erosion under carryover that
Section 3 measures.

\subsection{B.6 What the reduction removes, and what it does not}

Relative to the three-component matrix, the reduction removes:

\begin{itemize}
\item one within-component block, $A_{tv}$ ($n \times n$);
\item two cross-component blocks, $C_{tv,pb}$ and $C_{tv,br}$;
\item $n$ rows and columns from $R$, taking it from $2+3n$ to
  $2+2n$: at $n = 8$, from $26 \times 26$ to $18 \times 18$.
\end{itemize}

It removes no entry of \(b\), because the biomarker never coupled to
\(TV\). It removes no parameter from the architecture definitions of
Section 2.2, because neither references \(TV\).

One distinction is worth recording, because it is easy to conflate
two operations that are not equivalent. Removing the component is not
the same as flattening its trajectory. The code sets component
standard deviations independently of the Gompertz maximum, so setting
that maximum to zero yields a flat \(TV\) mean while leaving the \(TV\)
variance contribution, its block \(A_{tv}\), and its cross-component
blocks in place. A study wanting the simpler mental picture without
the change in residual variance should flatten; a study wanting the
smaller matrix should remove. This paper removes.

\bigskip\noindent

\rule{\textwidth}{0.4pt}\bigskip

\section{Appendix C. Positive definiteness and the feasible range of
$c_{bm}$}

The correlation matrix of Appendix B must be positive definite for
the multivariate normal draw to be valid. It is not automatically so:
the moderation parameter \(c_{bm}\) is bounded above, and the bound
depends on the correlation structure, the trial design, and the
number of occasions. This appendix derives the bound, gives a closed
form for the compound-symmetric case, and reports where the reduction
of Appendix B does and does not relax it.

The motivation is practical. Simulation drivers in this program
repair a non-positive-definite matrix silently, by projection onto
the nearest positive-definite matrix, so a grid cell that violates
the bound does not fail. It quietly becomes a different
data-generating process, with a realized biomarker-response
correlation no longer equal to the nominal \(c_{bm}\). Knowing the
bound in advance is therefore a precondition for trusting a grid.

\subsection{C.1 The constraint is a Schur complement on the biomarker
row}

Partition the correlation matrix of Appendix B.1 by separating the
biomarker from everything else. Because \(\mathrm{BL}\) is uncorrelated
with all other variables, it contributes an identity block that
factors out of every determinant and may be dropped. What remains is

\[
R' =
\begin{pmatrix}
1 & \tilde b^{\top} \\
\tilde b & M
\end{pmatrix},
\qquad
M =
\begin{pmatrix}
A_{pb} & C_{pb,br}\\
C_{br,pb} & A_{br}
\end{pmatrix},
\]

where \(M\) is the \(Kn \times Kn\) component block (\(K = 2\) here, \(K =
3\) under the full decomposition) and \(\tilde b \in \mathbb{R}^{Kn}\)
is the biomarker coupling extended by zeros: it equals \(b\) on the
\(BR\) positions and vanishes elsewhere.

By the Schur complement, \(R'\) is positive definite if and only if \(M\)
is positive definite and

\[
1 - \tilde b^{\top} M^{-1} \tilde b > 0 .
\]

Writing \(\tilde b = c_{bm} v\), where \(v\) is the unit-scaled coupling
pattern with \(v_p = \phi(t_{sd,p})\) on the \(BR\) positions and zero
elsewhere, the condition becomes \(c_{bm}^{2}\, v^{\top} M^{-1} v <
1\), giving the exact ceiling

\begin{equation}
c_{bm}^{\ast} \;=\; \bigl(v^{\top} M^{-1} v\bigr)^{-1/2}.
\label{eq:ceiling}
\end{equation}

Three consequences follow immediately, and they organize the rest of
this appendix. The ceiling depends on the component block only
through the quadratic form \(v^{\top} M^{-1} v\), so components that
are weakly coupled to \(BR\) contribute little to it. It depends on the
design through \(v\), which encodes the on-drug and off-drug pattern.
And it is unaffected by mean moderation, under which \(b \equiv 0\) and
the constraint reduces to positive definiteness of \(M\) alone.

\subsection{C.2 Closed form for the component block under compound
symmetry}

When the within-component blocks are compound symmetric rather than
AR(1), \(M\) admits an exact spectral decomposition. This is the
structure of the published implementation discussed in paper 01, and
it is worth recording because it explains a result of C.3 that is
otherwise surprising.

Under compound symmetry with a common within-component correlation
\(\rho\), common same-occasion cross-component correlation \(c_1\), and
common different-occasion cross-component correlation \(c_{\times}\),
the component block is a Kronecker sum,

\begin{equation}
M \;=\; A \otimes I_n \;+\; B \otimes (J_n - I_n),
\label{eq:kron}
\end{equation}

where \(J_n\) is the \(n \times n\) matrix of ones and \(A\), \(B\) are
\(K \times K\):

\[
A = I_K + c_1 (J_K - I_K), \qquad
B = \rho I_K + c_{\times}(J_K - I_K).
\]

\(A\) describes the same-occasion structure across components, \(B\) the
across-occasion structure. Since \(J_n - I_n\) has eigenvalues \(n-1\)
once and \(-1\) with multiplicity \(n-1\), and \(I_n\) commutes with it,
the \(Kn\) eigenvalues of \(M\) collapse to those of two \(K \times K\)
reductions:

\[
\mathrm{eig}(M) \;=\;
\underbrace{\mathrm{eig}\bigl(A + (n-1)B\bigr)}_{\text{once each}}
\;\cup\;
\underbrace{\mathrm{eig}\bigl(A - B\bigr)}_{\text{multiplicity } n-1}.
\]

Both reductions are of the form \(\alpha I_K + \beta (J_K - I_K)\),
whose eigenvalues are \(\alpha + (K-1)\beta\) once and \(\alpha - \beta\)
with multiplicity \(K-1\). Substituting gives all four distinct
eigenvalues in closed form:

\[
\begin{aligned}
\text{all-ones time direction:}\quad
&\bigl[1 + (n-1)\rho\bigr] + (K-1)\bigl[c_1 + (n-1)c_{\times}\bigr],\\
&\bigl[1 + (n-1)\rho\bigr] - \bigl[c_1 + (n-1)c_{\times}\bigr]
  \quad (K-1 \text{ times}),\\[4pt]
\text{time-contrast directions:}\quad
&(1-\rho) + (K-1)(c_1 - c_{\times}),\\
&(1-\rho) - (c_1 - c_{\times})
  \quad (K-1 \text{ times, each with multiplicity } n-1).
\end{aligned}
\]

At the program's reference values \(\rho = 0.7\), \(c_1 = 0.2\),
\(c_{\times} = 0.1\), \(n = 8\), the three-component block has exactly
four distinct eigenvalues, \(7.7\), \(5.0\), \(0.5\) and \(0.2\), the last
with multiplicity \(2(n-1) = 14\).

\subsection{C.3 Why the component count barely matters}

The smallest eigenvalue in C.2 is generally the last,

\begin{equation}
\lambda_{\min}(M) = (1-\rho) - (c_1 - c_{\times}),
\label{eq:lmin}
\end{equation}

and the striking feature of \eqref{eq:lmin} is what does not appear
in it. It contains neither \(n\) nor \(K\). Under compound symmetry the
positive-definiteness of the component block is independent of both
the number of occasions and the number of components. Adding a third
component does not make \(M\) harder to keep positive definite; nor
does adding occasions.

This is the analytic explanation for a result that would otherwise be
puzzling. One might expect that deleting a third of the rows and
columns from a correlation matrix would materially enlarge the
feasible parameter space. It does not, and \eqref{eq:lmin} shows why:
the binding structural constraint never depended on the count in the
first place.

Two further readings of \eqref{eq:lmin} are worth recording. The
operative cross-component quantity is the \emph{difference}
\(c_1 - c_{\times}\), not either coefficient separately, so keeping the
same-occasion and across-occasion couplings close to one another
matters more than keeping either small. And positive definiteness
requires \(\rho < 1 - (c_1 - c_{\times})\), which at the reference
values is \(\rho < 0.9\): a compound-symmetric structure cannot
represent very high autocorrelation at all.

Under AR(1) the corresponding statement is different, and no equally
compact closed form is available, because \(A_c\) is Toeplitz rather
than exchangeable and does not commute with the cross-component
structure in the way \eqref{eq:kron} requires. Numerically, the AR(1)
component block becomes progressively more ill-conditioned as
occasions are added, since each new occasion is strongly correlated
with its neighbours. The number of \emph{occasions} therefore does
constrain \(c_{bm}\) under AR(1), while the number of \emph{components}
still does not.

\subsection{C.4 Numerical ceilings}

Evaluating \eqref{eq:ceiling} directly at the program's reference
parameters (\(\rho = 0.7\), \(c_1 = 0.2\), \(c_{\times} = 0.1\), \(n = 8\),
AR(1) in cumulative weeks), minimized across each design's
randomization paths:

{\def\LTcaptype{none} % do not increment counter
\begin{longtable}[]{@{}llll@{}}
\toprule\noalign{}
Design & \(t_{1/2}\) & Three components & Two components \\
\midrule\noalign{}
\endhead
\bottomrule\noalign{}
\endlastfoot
CO & 0.0 & 0.619 & 0.627 \\
CO & 0.5 & 0.619 & 0.627 \\
CO & 1.0 & 0.619 & 0.627 \\
Hybrid & 0.0 & 0.456 & 0.469 \\
Hybrid & 0.5 & 0.506 & 0.516 \\
Hybrid & 1.0 & 0.544 & 0.552 \\
OL+BDC & 0.0 & 0.452 & 0.464 \\
OL+BDC & 0.5 & 0.500 & 0.509 \\
OL+BDC & 1.0 & 0.533 & 0.540 \\
\end{longtable}
}

The reduction buys between \(0.008\) and \(0.013\) of feasible range.
This is consistent with C.3: the deleted blocks were coupled to \(BR\)
only through \(c_1\) and \(c_{\times}\), both small, so they contributed
little to the quadratic form \(v^{\top} M^{-1} v\) that sets the
ceiling.

Three further observations follow from evaluating \eqref{eq:ceiling}
under variations of the structure, none of which is specific to the
reduction.

First, the ceiling rises with the carryover half-life within the two
non-crossover designs, from \(0.456\) to \(0.544\) under Hybrid. A longer
half-life makes the coupling vector \(v\) smoother, and a smoother \(v\)
aligns better with the well-conditioned directions of \(M\). The
crossover design is insensitive to \(t_{1/2}\) because its off-drug
occasions are spaced at least 2.5 weeks apart, so \(\phi(t_{sd})\) has
decayed to near zero at every one of them regardless.

Second, the number of occasions matters considerably more than the
number of components. Holding the design shape fixed and varying \(n\)
alone, the ceiling falls from approximately \(0.81\) at \(n = 4\) to
\(0.58\) at \(n = 12\) under a three-component structure. That is an
order of magnitude larger than the effect of deleting an entire
component.

Third, the exponent scale in \(A_c\) is consequential. Using the
occasion-index difference \(|i-j|\) in place of the cumulative week gap
\(|w_i - w_j|\), an otherwise identical AR(1) structure, lowers the
worst-cell ceiling from approximately \(0.45\) to \(0.35\). The code uses
the week gap; documentation elsewhere in this project has described
it as the index difference, and the two are not interchangeable.

\subsection{C.5 Practical implications}

For the present paper, \eqref{eq:ceiling} confirms that the grid of
Section 2.4 is feasible throughout: the largest moderation level used
is \(c_{bm} = 0.45\), against a worst-cell ceiling of \(0.464\) under the
two-component structure. The margin is real but not generous, and a
study wishing to extend the grid past \(c_{bm} = 0.45\) should evaluate
\eqref{eq:ceiling} first rather than rely on the simulation to
complain, since it will not.

More generally, the appendix identifies which levers move the
feasible range and which do not. The number of latent components does
not, and neither does the choice between two and three of them. The
number of measurement occasions does, substantially. The
within-component correlation \(\rho\) does, non-monotonically, with the
reference value \(0.7\) close to the maximum of the ceiling. And the
cross-component coefficients matter through their difference rather
than their magnitudes.

\paragraph{Provenance.}

The Kronecker identity \eqref{eq:kron} and
the eigenvalue decomposition of C.2 were verified numerically against
direct construction for \(K \in \{2,3\}\) and \(n \in \{4,8,12\}\). The
ceilings in C.4 were computed from \eqref{eq:ceiling} and agree with
an independent grid search over \(c_{bm}\) in steps of \(0.01\) to within
the grid step. The observations in C.4 concerning occasion count,
exponent scale, and \(\rho\) were obtained by the same grid search
under a reimplemented matrix builder rather than the production
\texttt{buildSigma}, and should be read as characterizing the
structure rather than as production audit results.

\section{References}\label{references}

\protect\phantomsection\label{refs}
\begin{CSLReferences}{0}{1}
\bibitem[\citeproctext]{ref-hendrickson2020}
\CSLLeftMargin{1. }%
\CSLRightInline{Hendrickson RC, Thomas RG, Schork NJ, Raskind MA. Optimizing aggregated {N}-of-1 trial designs for predictive biomarker validation: Statistical methods and theoretical findings. \emph{Frontiers in Digital Health}. 2020;2:13. doi:\href{https://doi.org/10.3389/fdgth.2020.00013}{10.3389/fdgth.2020.00013}}

\bibitem[\citeproctext]{ref-pmsimstats-paper11}
\CSLLeftMargin{2. }%
\CSLRightInline{{pmsimstats team}. A combined data-generating architecture for biomarker-treatment interactions: Channel super-additivity and mixing-weight sensitivity in aggregated {N}-of-1 trials. Published online 2026.}

\bibitem[\citeproctext]{ref-senn2019samplesize}
\CSLLeftMargin{3. }%
\CSLRightInline{Senn S. Sample size considerations for n-of-1 trials. \emph{Statistical Methods in Medical Research}. 2019;28(2):372-383. doi:\href{https://doi.org/10.1177/0962280217726801}{10.1177/0962280217726801}}

\bibitem[\citeproctext]{ref-pmsimstats-paper03}
\CSLLeftMargin{4. }%
\CSLRightInline{{pmsimstats team}. Latent-class and mixture-model formulations for biomarker-treatment interaction in {N}-of-1 trials. Published online 2026.}

\bibitem[\citeproctext]{ref-haller2018}
\CSLLeftMargin{5. }%
\CSLRightInline{Haller B, Ulm K. A simulation study on estimating biomarker-treatment interaction effects in randomized trials with prognostic variables. \emph{Trials}. 2018;19(1):128.}

\bibitem[\citeproctext]{ref-simon2010}
\CSLLeftMargin{6. }%
\CSLRightInline{Simon R. Clinical trial designs for evaluating the medical utility of prognostic and predictive biomarkers in oncology. \emph{Personalized Medicine}. 2010;7(1):33-47.}

\bibitem[\citeproctext]{ref-freidlin2014}
\CSLLeftMargin{7. }%
\CSLRightInline{Freidlin B, Korn EL. Biomarker enrichment strategies: Matching trial design to biomarker credentials. \emph{Nature Reviews Clinical Oncology}. 2014;11(2):81-90.}

\bibitem[\citeproctext]{ref-grenet2020}
\CSLLeftMargin{8. }%
\CSLRightInline{Grenet G, Blanc C, Bardel C, Gueyffier F, Roy P. Comparison of crossover and parallel-group designs for the identification of a binary predictive biomarker of the treatment effect. \emph{Basic and Clinical Pharmacology and Toxicology}. 2020;126(1):59-64. doi:\href{https://doi.org/10.1111/bcpt.13293}{10.1111/bcpt.13293}}

\bibitem[\citeproctext]{ref-senn2016mastering}
\CSLLeftMargin{9. }%
\CSLRightInline{Senn S. Mastering variation: Variance components and personalised medicine. \emph{Statistics in Medicine}. 2016;35(7):966-977. doi:\href{https://doi.org/10.1002/sim.6739}{10.1002/sim.6739}}

\bibitem[\citeproctext]{ref-wang2019}
\CSLLeftMargin{10. }%
\CSLRightInline{Wang Y, Schork NJ. Power and design issues in crossover-based {N}-of-1 clinical trials with fixed data collection periods. \emph{Healthcare (Basel)}. 2019;7(3):84.}

\end{CSLReferences}

\end{document}
